\documentclass[main.tex]{subfiles}


\begin{document}

%from pagenumber to up
% \phantomsection 
% \hypertarget{toc}{}

 

 

% номер секции вручную
\setcounter{section}{29
}
\addtocounter{section}{-1} 

\section{Индукция и линии магнитного поля}


\textbf{Индукция} ($\vec B$ [Тл])~--- это характеристика силовой способности магнитного поля в точке пространства. В каждой точке пространства вектор~$\vec B$ направлен туда же, куда и северный конец ($N$"=полюс) \emph{магнитной стрелки}, помещенной в данную точку (рис.\,\ref{fig:p140}, слева). 


    \begin{figure}[H]
    \centering

    \begin{tikzpicture}[font=\small,            line cap=round,                             >={Stealth[bend,inset=0pt,round,length=3mm, angle'=20]},vec/.style={>={Stealth[inset=0pt,round,length=3.5mm, angle'=20]}}]


\tikzset{ne/.pic={
  \path[rotate=0] (0,0) ++ (0,\dc) coordinate (ac) --++ ({3*\dc},{-\dc}) coordinate (rc) --++ ({-3*\dc},{-\dc}) coordinate (bc) --++ ({-3*\dc},{\dc}) coordinate (lc) --++ ({3*\dc},{\dc});
  \fill[red,semitransparent] (0,\dc) node[above,black,opaque] {6} --++ ({3*\dc},{-\dc}) --++ ({-3*\dc},{-\dc}) -- cycle;
  \fill[NavyBlue,semitransparent] (0,-\dc)  --++ ({-3*\dc},{\dc}) --++ ({3*\dc},{\dc}) -- cycle;
  \draw[black] (0,\dc)  --++ ({3*\dc},{-\dc}) --++ ({-3*\dc},{-\dc}) --++ ({-3*\dc},{\dc}) -- cycle;}
}


\def\lm{1} 
\def\dc{.1}


\begin{scope}[even odd rule,gray]
\def\dg{1.1} % должно сходиться с \dg в %3 линии
\clip[] ({-\dg*1cm-7*\pgflinewidth},{{-2.52*1cm-9*\pgflinewidth}}) rectangle++ ({2*\lm*1cm+\dg*1cm+\dg*1cm+14*\pgflinewidth},{2*2.52*1cm+4*\pgflinewidth})
(0,{-\lm/4}) rectangle++ ({2*\lm},{\lm/2});
%lines
%1
\def\ds{.2}
% \draw[yshift=.19cm,->,decoration={markings,mark=at position {1/2+3/20} with {\pic {ne};}},postaction={decorate}] ({2*\lm*1cm+\ds*1cm},{(\lm/4)*7/8}) 
% arc [start angle=0, end angle=90, x radius={1cm+\ds*1cm}, y radius=.5cm]
% coordinate (e1)
% ;

\path[yshift=.19cm,->] ({2*\lm*1cm+\ds*1cm},{(\lm/4)*7/8}) 
arc [start angle=0, end angle=90, x radius={1cm+\ds*1cm}, y radius=.5cm]
coordinate (e1)
[
    postaction=decorate,
    decoration={
      name=markings,
      mark={
        at position {1/2} with
        \fill[NavyBlue,semitransparent] (0,\dc) --++ ({3*\dc},{-\dc}) --++ ({-3*\dc},{-\dc}) node[above,black,opaque] {3} -- cycle;
        \fill[red,semitransparent] (0,-\dc)  --++ ({-3*\dc},{\dc}) --++ ({3*\dc},{\dc}) -- cycle;
        \draw[black] (0,\dc)  --++ ({3*\dc},{-\dc}) --++ ({-3*\dc},{-\dc}) --++ ({-3*\dc},{\dc}) -- cycle;
      }
    }
  ]
;

\path[yshift=.19cm] (e1) 
arc [start angle=90, end angle=180, x radius={1cm+\ds*1cm}, y radius=.5cm] 
arc [start angle=180, end angle=360, x radius={1cm+\ds*1cm}, y radius=.4cm] 
;


%2
\def\df{.5}
\path[yshift=.42cm,->] ({2*\lm*1cm+\df*1cm},{(\lm/4)*7/8}) 
arc [start angle=0, end angle=90, x radius={1cm+\df*1cm}, y radius=.9cm] 
coordinate (e1)
;

\path[yshift=.42cm] (e1) 
arc [start angle=90, end angle=180, x radius={1cm+\df*1cm}, y radius=.9cm] 
arc [start angle=180, end angle=360, x radius={1cm+\df*1cm}, y radius=.7cm] 
;


%3
\def\dg{1.1}
\path[yshift=.8cm,->] ({2*\lm*1cm+\dg*1cm},{(\lm/4)*7/8}) 
arc [start angle=0, end angle=90, x radius={1cm+\dg*1cm}, y radius=1.5cm] 
coordinate (e1)
[
    postaction=decorate,
    decoration={
      name=markings,
      mark={
        at position {.07} with
        \fill[NavyBlue,semitransparent] (0,\dc) node[left,black,opaque] {2}  --++ ({3*\dc},{-\dc}) --++ ({-3*\dc},{-\dc}) -- cycle;
        \fill[red,semitransparent] (0,-\dc)  --++ ({-3*\dc},{\dc}) --++ ({3*\dc},{\dc}) -- cycle;
        \draw[black] (0,\dc)  --++ ({3*\dc},{-\dc}) --++ ({-3*\dc},{-\dc}) --++ ({-3*\dc},{\dc}) -- cycle;
      }
    }
  ]
;

\path[yshift=.8cm] (e1) 
arc [start angle=90, end angle=180, x radius={1cm+\dg*1cm}, y radius=1.5cm] 
arc [start angle=180, end angle=360, x radius={1cm+\dg*1cm}, y radius=1.1cm] 
[
    postaction=decorate,
    decoration={
      name=markings,
      mark={
        at position {1/7} with
        \fill[NavyBlue,semitransparent] (0,\dc) node[below,black,opaque] {4} --++ ({3*\dc},{-\dc}) --++ ({-3*\dc},{-\dc}) -- cycle;
        \fill[red,semitransparent] (0,-\dc)  --++ ({-3*\dc},{\dc}) --++ ({3*\dc},{\dc}) -- cycle;
        \draw[black] (0,\dc)  --++ ({3*\dc},{-\dc}) --++ ({-3*\dc},{-\dc}) --++ ({-3*\dc},{\dc}) -- cycle;
      }
    }
  ]
;


%%%%%
%%%%%


%reverslines
\begin{scope}[yscale=-1]
%1
\def\ds{.2}
\path[yshift=.19cm,->] ({2*\lm*1cm+\ds*1cm},{(\lm/4)*7/8}) 
arc [start angle=0, end angle=90, x radius={1cm+\ds*1cm}, y radius=.5cm]
coordinate (e1)
;

\path[yshift=.19cm] (e1) 
arc [start angle=90, end angle=180, x radius={1cm+\ds*1cm}, y radius=.5cm] 
arc [start angle=180, end angle=360, x radius={1cm+\ds*1cm}, y radius=.4cm] 
;


%2
\def\df{.5}
\path[yshift=.42cm,->] ({2*\lm*1cm+\df*1cm},{(\lm/4)*7/8}) 
arc [start angle=0, end angle=90, x radius={1cm+\df*1cm}, y radius=.9cm] 
coordinate (e1)
;

\path[yshift=.42cm] (e1) 
arc [start angle=90, end angle=180, x radius={1cm+\df*1cm}, y radius=.9cm] 
arc [start angle=180, end angle=360, x radius={1cm+\df*1cm}, y radius=.7cm] 
[
    postaction=decorate,
    decoration={
      name=markings,
      mark={
        at position {1/6} with
        \fill[NavyBlue,semitransparent] (0,\dc)  --++ ({3*\dc},{-\dc}) --++ ({-3*\dc},{-\dc}) node[above,black,opaque] {5} -- cycle;
        \fill[red,semitransparent] (0,-\dc)  --++ ({-3*\dc},{\dc}) --++ ({3*\dc},{\dc}) -- cycle;
        \draw[black] (0,\dc)  --++ ({3*\dc},{-\dc}) --++ ({-3*\dc},{-\dc}) --++ ({-3*\dc},{\dc}) -- cycle;
      }
    }
  ]
;

%3
\def\dg{1.1}
\path[yshift=.8cm,->] ({2*\lm*1cm+\dg*1cm},{(\lm/4)*7/8}) 
arc [start angle=0, end angle=90, x radius={1cm+\dg*1cm}, y radius=1.5cm] 
coordinate (e1) pic {ne}
;

\path[yshift=.8cm] (e1) 
arc [start angle=90, end angle=180, x radius={1cm+\dg*1cm}, y radius=1.5cm] 
arc [start angle=180, end angle=360, x radius={1cm+\dg*1cm}, y radius=1.1cm] 
;
\end{scope}
\end{scope}





%magnet
\fill[red,semitransparent] (0,{-\lm/4}) rectangle++ (\lm,{\lm/2});
\fill[NavyBlue,semitransparent] (\lm,{-\lm/4}) rectangle++ (\lm,{\lm/2});

\node[right] at (0,.) {$S$}; 
\node[left] at (2,.) {$N$}; 

\draw (0,{-\lm/4}) rectangle++ ({2*\lm},{\lm/2});

% \path[gray] (0,0) --++ (-1,0)
% [
%     postaction=decorate,
%     decoration={
%       name=markings,
%       mark={
%         at position {1} with
%         \fill[NavyBlue,semitransparent] (0,\dc)  --++ ({3*\dc},{-\dc}) --++ ({-3*\dc},{-\dc}) -- cycle;
%         \fill[red,semitransparent] (0,-\dc)  --++ ({-3*\dc},{\dc}) --++ ({3*\dc},{\dc}) -- cycle;
%         \draw[black] (0,\dc)  --++ ({3*\dc},{-\dc}) --++ ({-3*\dc},{-\dc}) --++ ({-3*\dc},{\dc}) -- cycle;
%       }
%     }
%   ]
% ;

\path[gray] ({2*\lm},0) --++ (1,0)
[
    postaction=decorate,
    decoration={
      name=markings,
      mark={
        at position {1/2} with
        \fill[NavyBlue,semitransparent] (0,\dc) node[above,black,opaque] {1} --++ ({3*\dc},{-\dc}) --++ ({-3*\dc},{-\dc}) -- cycle;
        \fill[red,semitransparent] (0,-\dc)  --++ ({-3*\dc},{\dc}) --++ ({3*\dc},{\dc}) -- cycle;
        \draw[black] (0,\dc)  --++ ({3*\dc},{-\dc}) --++ ({-3*\dc},{-\dc}) --++ ({-3*\dc},{\dc}) -- cycle;
      }
    }
  ]
;


%phantom
\phantom{
\def\mas{1.5}
\path[gray] (0,0) --++ (-1,0)
[
    postaction=decorate,
    decoration={
      name=markings,
      mark={
        at position {1/2} with
        \draw[gray,vec,thick,->] (0,0)  --++ ({\mas*1},0) node[below,black] {$\vec B_1$};
      }
    }
  ]
;
}







% \draw[tips, -{>[length=10pt,bend]}] (0,0) to[bend left] (1,0);












\def\mas{1.5}

\begin{scope}[xshift=8cm]

\begin{scope}[even odd rule,gray]
\def\dg{1.1} % должно сходиться с \dg в %3 линии
\clip[] ({-\dg*1cm-7*\pgflinewidth},{{-2.52*1cm-9*\pgflinewidth}}) rectangle++ ({2*\lm*1cm+\dg*1cm+\dg*1cm+14*\pgflinewidth},{2*2.52*1cm+4*\pgflinewidth})
(0,{-\lm/4}) rectangle++ ({2*\lm},{\lm/2});
%lines
%1
\def\ds{.2}
% \draw[yshift=.19cm,->,decoration={markings,mark=at position {1/2+3/20} with {\pic {ne};}},postaction={decorate}] ({2*\lm*1cm+\ds*1cm},{(\lm/4)*7/8}) 
% arc [start angle=0, end angle=90, x radius={1cm+\ds*1cm}, y radius=.5cm]
% coordinate (e1)
% ;

\path[yshift=.19cm,->] ({2*\lm*1cm+\ds*1cm},{(\lm/4)*7/8}) 
arc [start angle=0, end angle=90, x radius={1cm+\ds*1cm}, y radius=.5cm]
coordinate (e1)
[
    postaction=decorate,
    decoration={
      name=markings,
      mark={
        at position {1/2} with
        \draw[vec,thick,->] (0,0)  --++ ({\mas*.8},0) node[above,black] {$\vec B_3$};
        \node at (0,0) [circle,fill=black,inner sep=0pt,minimum size=1mm,label=below:] {};
      }
    }
  ]
;

\path[yshift=.19cm] (e1) 
arc [start angle=90, end angle=180, x radius={1cm+\ds*1cm}, y radius=.5cm] 
arc [start angle=180, end angle=360, x radius={1cm+\ds*1cm}, y radius=.4cm] 
;


%2
\def\df{.5}
\path[yshift=.42cm,->] ({2*\lm*1cm+\df*1cm},{(\lm/4)*7/8}) 
arc [start angle=0, end angle=90, x radius={1cm+\df*1cm}, y radius=.9cm] 
coordinate (e1)
;

\path[yshift=.42cm] (e1) 
arc [start angle=90, end angle=180, x radius={1cm+\df*1cm}, y radius=.9cm] 
arc [start angle=180, end angle=360, x radius={1cm+\df*1cm}, y radius=.7cm] 
;


%3
\def\dg{1.1}
\path[yshift=.8cm,->] ({2*\lm*1cm+\dg*1cm},{(\lm/4)*7/8}) 
arc [start angle=0, end angle=90, x radius={1cm+\dg*1cm}, y radius=1.5cm] 
coordinate (e1)
[
    postaction=decorate,
    decoration={
      name=markings,
      mark={
        at position {.07} with
        \draw[vec,thick,->] (0,0)  --++ ({\mas*.7},0) node[yshift=-.03cm,left,black] {$\vec B_2$};
        \node at (0,0) [circle,fill=black,inner sep=0pt,minimum size=1mm,label=below:] {};
      }
    }
  ]
;

\path[yshift=.8cm] (e1) 
arc [start angle=90, end angle=180, x radius={1cm+\dg*1cm}, y radius=1.5cm] 
arc [start angle=180, end angle=360, x radius={1cm+\dg*1cm}, y radius=1.1cm] 
[
    postaction=decorate,
    decoration={
      name=markings,
      mark={
        at position {1/7} with
        \draw[vec,thick,->] (0,0)  --++ ({\mas*.5},0) node[below,black] {$\vec B_4$};
        \node at (0,0) [circle,fill=black,inner sep=0pt,minimum size=1mm,label=below:] {};
      }
    }
  ]
;


%%%%%
%%%%%


%reverslines
\begin{scope}[yscale=-1]
%1
\def\ds{.2}
\path[yshift=.19cm,->] ({2*\lm*1cm+\ds*1cm},{(\lm/4)*7/8}) 
arc [start angle=0, end angle=90, x radius={1cm+\ds*1cm}, y radius=.5cm]
coordinate (e1)
;

\path[yshift=.19cm] (e1) 
arc [start angle=90, end angle=180, x radius={1cm+\ds*1cm}, y radius=.5cm] 
arc [start angle=180, end angle=360, x radius={1cm+\ds*1cm}, y radius=.4cm] 
;


%2
\def\df{.5}
\path[yshift=.42cm,->] ({2*\lm*1cm+\df*1cm},{(\lm/4)*7/8}) 
arc [start angle=0, end angle=90, x radius={1cm+\df*1cm}, y radius=.9cm] 
coordinate (e1)
;

\path[yshift=.42cm] (e1) 
arc [start angle=90, end angle=180, x radius={1cm+\df*1cm}, y radius=.9cm] 
arc [start angle=180, end angle=360, x radius={1cm+\df*1cm}, y radius=.7cm] 
[
    postaction=decorate,
    decoration={
      name=markings,
      mark={
        at position {1/6} with
        \draw[gray,vec,thick,->] (0,0)  --++ ({\mas*.6},0) node[above,black] {$\vec B_5$};
        \node at (0,0) [circle,fill=black,inner sep=0pt,minimum size=1mm,label=below:] {};
      }
    }
  ]
;

%3
\def\dg{1.1}
\path[yshift=.8cm,->] ({2*\lm*1cm+\dg*1cm},{(\lm/4)*7/8}) 
arc [start angle=0, end angle=90, x radius={1cm+\dg*1cm}, y radius=1.5cm] 
coordinate (e1)
;
\draw[vec,thick,->] (e1)  --++ ({\mas*-.4},0) node[above,black] {$\vec B_6$};
\node at (e1) [circle,fill=black,inner sep=0pt,minimum size=1mm,label=below:] {};


\path[yshift=.8cm] (e1) 
arc [start angle=90, end angle=180, x radius={1cm+\dg*1cm}, y radius=1.5cm] 
arc [start angle=180, end angle=360, x radius={1cm+\dg*1cm}, y radius=1.1cm] 
;
\end{scope}
\end{scope}





%magnet
\fill[red,semitransparent] (0,{-\lm/4}) rectangle++ (\lm,{\lm/2});
\fill[NavyBlue,semitransparent] (\lm,{-\lm/4}) rectangle++ (\lm,{\lm/2});

\node[right] at (0,.) {$S$}; 
\node[left] at (2,.) {$N$}; 

\draw (0,{-\lm/4}) rectangle++ ({2*\lm},{\lm/2});

% \path[gray] (0,0) --++ (-1,0)
% [
%     postaction=decorate,
%     decoration={
%       name=markings,
%       mark={
%         at position {1} with
%         \draw[vec,thick,->] (0,0) --++ (-.8,0);
%       }
%     }
%   ]
% ;

\path[gray] ({2*\lm},0) --++ (1,0)
[
    postaction=decorate,
    decoration={
      name=markings,
      mark={
        at position {1/2} with
        \draw[gray,vec,thick,->] (0,0)  --++ ({\mas*1},0) node[below,black] {$\vec B_1$};
        \node at (0,0) [circle,fill=black,inner sep=0pt,minimum size=1mm,label=below:] {};
      }
    }
  ]
;








% \draw[tips, -{>[length=10pt,bend]}] (0,0) to[bend left] (1,0);

    
\end{scope}


    \end{tikzpicture}
\caption{Стрелки и соответствующие векторы индукции}
\label{fig:p140}
\end{figure}


%%%%%%%%%%%%%%%%%%%%%%%%%%%%%%%%%%%%%
%%%%%%%%%%%%%%%%%%%%%%%%%%%%%%%%%%%%%
%%%%%%%%%%%%%%%%%%%%%%%%%%%%%%%%%%%%%

Возле магнита расположены шесть магнитных стрелок, способных свободно вращаться вокруг своих центров. На рис.\,\ref{fig:p140} (справа) изображены шесть векторов индукции поля магнита: каждой магнитной стрелке отвечает свой вектор~$\vec B$.  

Наглядно описывать магнитное поле в пространстве принято с помощью \textbf{линий магнитного поля}~--- линий, идущих вдоль векторов индукции поля в различных точках пространства\footnote{Эти линии называют также \emph{линиями индукции} или \emph{магнитными силовыми линиями}. Всякая линия поля изображается так, что в любой ее точке вектор индукции направлен по касательной к линии поля.}.

На рис.\,\ref{fig:p141} показана картина линий поля полосового магнита.


    \begin{figure}[H]
    \centering

    \begin{tikzpicture}[font=\small,            line cap=round,                             >={Stealth[bend,inset=0pt,round,length=3mm, angle'=20]},vec/.style={>={Stealth[inset=0pt,round,length=3.5mm, angle'=20]}}]


\def\lm{1} 

\begin{scope}[even odd rule,gray]
\def\dg{1.1} % должно сходиться с \dg в %3 линии
\clip[] ({-\dg*1cm-\pgflinewidth},{{-2.52*1cm-4*\pgflinewidth}}) rectangle++ ({2*\lm*1cm+\dg*1cm+\dg*1cm+2*\pgflinewidth},{2*2.52*1cm+8*\pgflinewidth})
(0,{-\lm/4}) rectangle++ ({2*\lm},{\lm/2});
;
%lines
%1
\def\ds{.2}
\draw[yshift=.19cm,->] ({2*\lm*1cm+\ds*1cm},{(\lm/4)*7/8}) 
arc [start angle=0, end angle=90, x radius={1cm+\ds*1cm}, y radius=.5cm]
coordinate (e1)
;

\draw[yshift=.19cm] (e1) 
arc [start angle=90, end angle=180, x radius={1cm+\ds*1cm}, y radius=.5cm] 
arc [start angle=180, end angle=360, x radius={1cm+\ds*1cm}, y radius=.4cm] 
;


%2
\def\df{.5}
\draw[yshift=.42cm,->] ({2*\lm*1cm+\df*1cm},{(\lm/4)*7/8}) 
arc [start angle=0, end angle=90, x radius={1cm+\df*1cm}, y radius=.9cm] 
coordinate (e1)
;

\draw[yshift=.42cm] (e1) 
arc [start angle=90, end angle=180, x radius={1cm+\df*1cm}, y radius=.9cm] 
arc [start angle=180, end angle=360, x radius={1cm+\df*1cm}, y radius=.7cm] 
;


%3
\def\dg{1.1}
\draw[yshift=.8cm,->] ({2*\lm*1cm+\dg*1cm},{(\lm/4)*7/8}) 
arc [start angle=0, end angle=90, x radius={1cm+\dg*1cm}, y radius=1.5cm] 
coordinate (e1)
;

\draw[yshift=.8cm] (e1) 
arc [start angle=90, end angle=180, x radius={1cm+\dg*1cm}, y radius=1.5cm] 
arc [start angle=180, end angle=360, x radius={1cm+\dg*1cm}, y radius=1.1cm] 
;


%%%%%
%%%%%


%reverslines
\begin{scope}[yscale=-1]
%1
\def\ds{.2}
\draw[yshift=.19cm,->] ({2*\lm*1cm+\ds*1cm},{(\lm/4)*7/8}) 
arc [start angle=0, end angle=90, x radius={1cm+\ds*1cm}, y radius=.5cm]
coordinate (e1)
;

\draw[yshift=.19cm] (e1) 
arc [start angle=90, end angle=180, x radius={1cm+\ds*1cm}, y radius=.5cm] 
arc [start angle=180, end angle=360, x radius={1cm+\ds*1cm}, y radius=.4cm] 
;


%2
\def\df{.5}
\draw[yshift=.42cm,->] ({2*\lm*1cm+\df*1cm},{(\lm/4)*7/8}) 
arc [start angle=0, end angle=90, x radius={1cm+\df*1cm}, y radius=.9cm] 
coordinate (e1)
;

\draw[yshift=.42cm] (e1) 
arc [start angle=90, end angle=180, x radius={1cm+\df*1cm}, y radius=.9cm] 
arc [start angle=180, end angle=360, x radius={1cm+\df*1cm}, y radius=.7cm] 
;


%3
\def\dg{1.1}
\draw[yshift=.8cm,->] ({2*\lm*1cm+\dg*1cm},{(\lm/4)*7/8}) 
arc [start angle=0, end angle=90, x radius={1cm+\dg*1cm}, y radius=1.5cm] 
coordinate (e1)
;

\draw[yshift=.8cm] (e1) 
arc [start angle=90, end angle=180, x radius={1cm+\dg*1cm}, y radius=1.5cm] 
arc [start angle=180, end angle=360, x radius={1cm+\dg*1cm}, y radius=1.1cm] 
;
\end{scope}
\end{scope}





%magnet
\fill[red,semitransparent] (0,{-\lm/4}) rectangle++ (\lm,{\lm/2});
\fill[NavyBlue,semitransparent] (\lm,{-\lm/4}) rectangle++ (\lm,{\lm/2});

\node[right] at (0,.) {$S$}; 
\node[left] at (2,.) {$N$}; 

\draw[gray,decoration={markings,mark=at position {1/2+1.5/20} with {\arrow{<}}},postaction={decorate}] (0,0) --++ (-2,0); 
\draw[gray,decoration={markings,mark=at position {1/2+1.5/20} with {\arrow{>}}},postaction={decorate}] ({2*\lm},0) --++ (2,0);

\draw (0,{-\lm/4}) rectangle++ ({2*\lm},{\lm/2});







% \draw[tips, -{>[length=10pt,bend]}] (0,0) to[bend left] (1,0);



    \end{tikzpicture}
\caption{Линии магнитного поля полосового магнита}
\label{fig:p141}
\end{figure}


\emph{Линии поля выходят из северного полюса магнита и входят в южный полюс (линии не пересекаются; внутри магнита они замыкаются). Чем больше густота линий, тем больше индукция поля в данной области пространства.}

% Так, в ситуации на рис.\,\ref{fig:p141}

Именно вдоль линий поля выстраиваются магнитные стрелки (или железные опилки), размещенные возле магнита. 










 
\end{document}